% ============================================================
% content/solutions.tex —— 2021 湖南卷物理 · 参考答案与详细解析
% 解析类型：source_embedded（转录自官方 Word 版内嵌解析，独立复核）
% ============================================================

\begin{answer}
\gksolhead{1}{单项选择题}
\gkanswer{答案 D}
半衰期是统计概念，指大量放射性元素的原子核衰变一半所用的时间，其表达式可写成
\[
  N=N_{0}\left(\frac{1}{2}\right)^{\frac{t}{\tau}},
\]
由此可知，放射性元素经过两个完整的半衰期后还剩四分之一，不会衰变殆尽，选项 A 错误。原子核衰变时核电荷数守恒，质量数守恒，选项 B 错误。放射性元素的半衰期与压力、温度、浓度等外部条件无关，选项 C 错误。过量放射性辐射对人体组织有破坏作用，辐射强度在安全剂量内则没有伤害，选项 D 正确。
\end{answer}

\begin{answer}
\gksolhead{2}{单项选择题}
\gkanswer{答案 D}
动量 $p=mv$，质点做初速度为零的匀加速直线运动，有
\[
  v^{2}=2ax,
\]
联立解得
\[
  p=m\sqrt{2ax},
\]
$p$ 随 $x$ 单调增大且曲线凸向 $x$ 轴（$p^{2}=2ma\cdot x$ 为开口向右、过原点的抛物线），由此可知选项 D 正确，A、B、C 错误。
\end{answer}

\begin{answer}
\gksolhead{3}{单项选择题}
\gkanswer{答案 C}
动车组匀加速启动过程中，根据牛顿第二定律有
\[
  F-kv=ma,
\]
因为加速度 $a$ 不变、速度 $v$ 改变，所以牵引力 $F$ 改变，选项 A 错误。由四节动力车厢输出功率均为额定值，可得
\[
  4P=Fv,\qquad F-kv=ma',
\]
联立解得 $a'=\dfrac{4P}{mv}-\dfrac{kv}{m}$，因为 $v$ 改变，所以 $a'$ 改变，选项 B 错误。由题干得
\[
  4P=F_{\text{阻}}v_{m}=kv_{m}^{2},
\]
若输出总功率为 $2.25P$，则有
\[
  2.25P=kv_{\text{匀}}^{2},
\]
联立解得 $v_{\text{匀}}=\dfrac{3}{4}v_{m}$，选项 C 正确。若四节动力车厢输出功率均为额定值，根据动能定理得
\[
  4Pt-W_{\text{阻}}=\frac{1}{2}mv_{m}^{2}-0,
\]
解得 $W_{\text{阻}}=4Pt-\dfrac{1}{2}mv_{m}^{2}$，选项 D 错误。
\end{answer}

\begin{answer}
\gksolhead{4}{单项选择题}
\gkanswer{答案 B}
如图所示，$-q$、$q$ 在 $P$ 点的合电场强度为
\[
  E_{1}=\sqrt{2}\,k\frac{q}{a^{2}},
\]
与水平方向夹角为 $45^{\circ}$。因为 $P$ 点电场强度为零，所以 $Q$ 必在 $E_{1}$ 的方向上，由几何关系得 $Q$ 的坐标为 $(0,2a)$，在 $P$ 点满足
\[
  \sqrt{2}\,k\frac{q}{a^{2}}=k\frac{Q}{(\sqrt{2}\,a)^{2}},
\]
解得 $Q=2\sqrt{2}\,q$，选项 B 正确。

\gkfig{figures/s4-e}
\end{answer}

\begin{answer}
\gksolhead{5}{单项选择题}
\gkanswer{答案 C}
如图所示，在任意位置，对小滑块受力分析，设 $F_{N}$ 与竖直方向的夹角为 $\theta$，根据平衡条件有
\[
  F_{N}=mg\cos\theta,\qquad F=mg\sin\theta.
\]
由 $A$ 向 $B$ 移动过程中，$\theta$ 角增大，所以推力 $F$ 增大、$F_{N}$ 减小，选项 A、B 错误。选小滑块和凹槽整体为研究对象，推力在水平方向的分力
\[
  F_{\text{水平}}=F\cos\theta=mg\sin\theta\cos\theta=\frac{mg\sin2\theta}{2},
\]
先增大后减小，所以墙面对凹槽的压力先增大后减小，选项 C 正确。推力在竖直方向的分力 $F_{\text{竖}}=F\sin\theta=mg\sin^{2}\theta$ 增大，故水平地面对凹槽的支持力 $F_{\text{支}}=(m+m')g-F_{\text{竖}}$ 减小，选项 D 错误。

\gkfig{figures/s5-force}
\end{answer}

\begin{answer}
\gksolhead{6}{单项选择题}
\gkanswer{答案 A}
在滑片从 $a$ 向 $b$ 滑动过程中，根据串并联电路特点可知，变压器负载电阻先增大后减小，滑片在中间时，电阻最大；将变压器等效成用电器，等效电阻为
\[
  R'=\frac{n_{1}^{2}}{n_{2}^{2}}\,R_{\text{负}},
\]
根据闭合电路欧姆定律可知，流过 $L_{1}$ 的电流 $I_{1}$ 先减小后增大，所以 $L_{1}$ 先变暗后变亮。滑片从 $a$ 向 $b$ 滑动过程中，滑到中点前，由于 $I_{1}$ 减小，$L_{1}$ 两端电压减小，变压器原线圈两端电压 $U_{1}$ 增大，副线圈两端电压 $U_{2}$ 增大，对于 $L_{2}$ 所在支路，电阻减小，所以电流增大，$L_{2}$ 变亮。滑片滑到中点后，由于 $I_{1}$ 增大，$U_{1}$ 将减小，$U_{2}$ 也减小，$R_{0}$ 所在支路电阻在增大，电流减小，而由于 $I_{1}$ 增大，副线圈中电流也增大，所以流过 $L_{2}$ 的电流增大，因此 $L_{2}$ 一直变亮，选项 A 正确。
\end{answer}

\begin{answer}
\gksolhead{7}{多项选择题}
\gkanswer{答案 AC}
根据万有引力定律 $F=G\dfrac{m_{\text{地}}m}{R^{2}}$，在地面上时 $F_{1}=\dfrac{Gm_{\text{地}}m}{R^{2}}$，进入轨道后
\[
  F_{2}=G\frac{m_{\text{地}}m}{\left(R+\frac{1}{16}R\right)^{2}},
\]
则
\[
  \frac{F_{2}}{F_{1}}=\frac{R^{2}}{\left(\frac{17}{16}\right)^{2}R^{2}}=\left(\frac{16}{17}\right)^{2},
\]
即 $F_{2}=\left(\dfrac{16}{17}\right)^{2}F_{1}$，选项 A 正确。由
\[
  \frac{Gm_{\text{地}}m}{r^{2}}=\frac{mv^{2}}{r}
\]
得 $v=\sqrt{\dfrac{Gm_{\text{地}}}{r}}$，因此核心舱在轨道上飞行的速度小于 \SI{7.9}{\kilo\metre\per\second}，选项 B 错误。由
\[
  \frac{Gm_{\text{地}}m}{r^{2}}=m\frac{4\pi^{2}}{T^{2}}r
\]
可得 $T=2\pi\sqrt{\dfrac{r^{3}}{Gm_{\text{地}}}}$，核心舱的高度小于地球同步卫星的高度，所以其飞行周期小于 \SI{24}{\hour}，选项 C 正确。加挂实验舱后，空间站质量增大，但轨道半径不变，选项 D 错误。
\end{answer}

\begin{answer}
\gksolhead{8}{多项选择题}
\gkanswer{答案 ABD}
0 到 $t_{1}$ 时间内，墙对 $B$ 的冲量与弹簧对 $B$ 的冲量大小相等，而弹簧对 $B$ 的冲量与弹簧对 $A$ 的冲量大小相等，根据动量定理，弹簧对 $A$ 的冲量 $I=m_{A}v_{0}$，所以选项 A 正确。分析运动过程可知，$t_{1}$ 时刻弹簧恢复原长，之后 $A$ 减速、$B$ 加速，当弹簧伸长量最大时二者速度相等，这时二者加速度均最大，即 $t_{2}$ 时刻，从图像可以看出，$B$ 的加速度大，所以质量小，即 $m_{B}<m_{A}$，选项 B 正确。初始时刻，系统机械能仅为弹性势能，弹簧形变量为 $x$；撤去外力后，系统机械能守恒，当 $B$ 运动后、$A$、$B$ 速度相等时，弹簧形变量最大，这时系统机械能包括动能和弹性势能，因此弹簧的最大形变量小于 $x$，选项 C 错误。$a$--$t$ 图像中图线与 $t$ 轴之间围的面积表示速度变化，设 $A$、$B$ 在 $t_{2}$ 时刻速度为 $v_{1}$，则
\[
  v_{0}=S_{1},\qquad v_{1}=v_{0}-S_{2},\qquad v_{1}=S_{3},
\]
解得 $S_{1}-S_{2}=S_{3}$，选项 D 正确。
\end{answer}

\begin{answer}
\gksolhead{9}{多项选择题}
\gkanswer{答案 AB}
由题意可得
\[
  U_{ab}=\varphi_{a}-\varphi_{b}=\frac{2W}{q},\qquad
  U_{cd}=\varphi_{c}-\varphi_{d}=\frac{W}{q}.
\]
过 $b$ 点作 $cd$ 的平行线，交圆于 $e$ 点，连接 $ce$，由几何关系可得四边形 $cObe$ 为平行四边形，故
\[
  \varphi_{c}-\varphi_{e}=\varphi_{O}-\varphi_{b}=\frac{W}{q},
\]
则 $\varphi_{c}-\varphi_{d}=\varphi_{c}-\varphi_{e}$，即 $\varphi_{d}=\varphi_{e}$，电场线沿 $ce$ 方向，即 $ab$ 方向，选项 A 正确。
\[
  W_{db}=\frac{U_{Ob}}{2}\cdot q=\frac{1}{2}\cdot\frac{W}{q}\cdot q=\frac{W}{2},
\]
选项 B 正确。沿电场线方向电势降低，$\varphi_{a}>\varphi_{c}$，选项 C 错误。从 $d$ 点射入的粒子，有可能沿电场线反方向做直线运动，选项 D 错误。
\end{answer}

\begin{answer}
\gksolhead{10}{多项选择题}
\gkanswer{答案 CD}
金属框匀速通过磁场时，重力与安培力相等，即
\[
  mg=BIl,\qquad I=\frac{Blv_{\bot}}{R},\qquad v_{\bot}^{2}=2gh,
\]
联立可得
\[
  B^{2}=\frac{mgR}{l^{2}\sqrt{2gh}},
\]
$B$ 与 $v_{0}$ 无关，但与 $h$ 的 $\dfrac{1}{4}$ 次方成反比，选项 A 错误。通过磁场过程中根据右手定则，金属框下边进入时电流为逆时针方向，金属框上边进入时电流为顺时针方向，选项 B 错误。因为金属框匀速通过磁场，重力做功与克服安培力做功相等，所以克服安培力做功的功率与重力做功的功率相等，选项 C 正确。因为通过磁场，重力做功不变，根据能量守恒定律得，产生的热量不变，选项 D 正确。
\end{answer}

\begin{answer}
\gksolhead{11}{实验题}
\gkanswer{答案 (1) $1.02$；(6) $0.346$}
(1) 该游标卡尺分度值为 \SI{0.1}{\milli\metre}，则
\[
  h=\SI{10.2}{\milli\metre}=\SI{1.02}{\centi\metre}.
\]

(6) 滑块沿倾斜气垫导轨下滑，由牛顿第二定律
\[
  mg\sin\theta=ma,
\]
其中
\[
  \sin\theta=\frac{nh}{l},
\]
解得
\[
  a=\frac{gh}{l}\cdot n.
\]
由所绘图像可得斜率 $k=\dfrac{0.519}{6}\,\si{\metre\per\second\squared}$，故
\[
  a_{4}=4k=\SI{0.346}{\metre\per\second\squared},
\]
也可直接由所绘图线读出第 4 组数据。所绘图像如图所示。

\gkfig{figures/s11-line}
\end{answer}

\begin{answer}
\gksolhead{12}{实验题}
\gkanswer{答案 (2) $b$；(4) $E=\dfrac{r_{0}}{k}$，$r=\dfrac{r_{0}d}{k}-R_{0}-R_{A}$；(5) 见解析}
(2) 开始时接入线路电阻应最大，金属夹应夹在电阻丝的 $b$ 端。

(4) 根据闭合电路欧姆定律得
\[
  E=I(R_{0}+R_{x}+R_{A}+r),
\]
其中 $R_{x}=r_{0}\theta$，解得
\[
  \frac{1}{I}=\frac{r_{0}}{E}\,\theta+\frac{R_{0}+R_{A}+r}{E}.
\]
由题图乙得
\[
  \frac{r_{0}}{E}=k,\qquad \frac{R_{0}+R_{A}+r}{E}=d,
\]
联立得
\[
  E=\frac{r_{0}}{k},\qquad r=\frac{r_{0}d}{k}-R_{0}-R_{A}.
\]

(5) 电路图如答案图所示。

\gkfig{figures/s12-circuit}

$S$ 先打到 $b$ 端，读出电流表读数 $I$；再把 $S$ 打到 $a$ 端，移动 $R_{x}$ 上的金属夹，使电流表的读数也为 $I$，此时 $R_{x}$ 的值与 $R_{0}$ 相等，记下偏转角度 $\theta$，则
\[
  r_{0}=\frac{R_{0}}{\theta}.
\]
\end{answer}

\begin{answer}
\gksolhead{13}{解答题}
\gkanswer{答案 (1) $B_{1}=\dfrac{mv}{qr_{1}}$；(2) $B_{2}=\dfrac{mv}{qr_{2}}$，垂直纸面向里，面积 $S=\pi r_{2}^{2}$}
\gkanswer{(3) $B_{3}=\dfrac{mv}{qr_{3}}$，$B_{4}=\dfrac{mv}{qr_{4}}$，$S_{\text{Ⅱ}}=\dfrac{1}{2}\pi r_{3}^{2}-r_{3}^{2}$，$S_{\text{Ⅳ}}=\dfrac{1}{2}\pi r_{4}^{2}-r_{4}^{2}$}
(1) 带电粒子在磁场中做匀速圆周运动，洛伦兹力提供向心力，有
\[
  qB_{1}v=m\frac{v^{2}}{R_{1}},
\]
粒子汇聚到 $O$ 点，故 $R_{1}=r_{1}$，得
\[
  B_{1}=\frac{mv}{qr_{1}}.
\]

(2) 磁场区域如图所示。

\gkphoto[0.55\textwidth]{photos/image22-field2.png}{第 (2) 问磁场区域解析图（原图保留，非矢量例外）}

同理有
\[
  qB_{2}v=m\frac{v^{2}}{R_{2}},
\]
粒子经磁场后沿 $x$ 轴正方向平行射出，故 $R_{2}=r_{2}$，得
\[
  B_{2}=\frac{mv}{qr_{2}},
\]
方向垂直纸面向里，磁场区域的面积 $S=\pi r_{2}^{2}$。

(3) 磁场区域如图所示。

\gkphoto[0.55\textwidth]{photos/image23-field3.png}{第 (3) 问磁场区域解析图（原图保留，非矢量例外）}

在区域 Ⅰ 中，有
\[
  qB_{3}v=m\frac{v^{2}}{R_{3}},
\]
且 $R_{3}=r_{3}$，得
\[
  B_{3}=\frac{mv}{qr_{3}}.
\]
在区域 Ⅲ 中，有
\[
  qB_{4}v=m\frac{v^{2}}{R_{4}},
\]
且 $R_{4}=r_{4}$，得
\[
  B_{4}=\frac{mv}{qr_{4}}.
\]
区域 Ⅱ 的面积为
\[
  S_{\text{Ⅱ}}=\frac{1}{2}\pi r_{3}^{2}-r_{3}^{2},
\]
区域 Ⅳ 的面积为
\[
  S_{\text{Ⅳ}}=\frac{1}{2}\pi r_{4}^{2}-r_{4}^{2}.
\]
\end{answer}

\begin{answer}
\gksolhead{14}{解答题}
\gkanswer{答案 (1) $v_{0}=\sqrt{2\mu gl}$；(2) $x^{2}+4y^{2}-8\mu ly=0$（$0<x\le2\mu l$）}
\gkanswer{(3) $\dfrac{3\lambda-1}{\lambda-3}\,\mu l<h\le\dfrac{4(\lambda^{2}+\lambda+1)}{(\lambda-1)^{2}}\,\mu l$}
(1) 根据动能定理得
\[
  mg\cdot 2\mu l-\mu mgl=\frac{1}{2}mv_{0}^{2}-0,
\]
解得
\[
  v_{0}=\sqrt{2\mu gl}.
\]

(2) 当 $A$ 落在 $P$ 点时
\[
  v_{0}'t_{1}=2\mu l,\qquad \frac{1}{2}gt_{1}^{2}=\mu l,
\]
\[
  E_{k}=\frac{1}{2}mv_{0}'^{2}+mg\mu l,
\]
联立解得 $E_{k}=2\mu mgl$。当 $A$ 落在任一点 $(x,y)$ 时
\[
  v_{0}''t_{2}=x,\qquad \frac{1}{2}gt_{2}^{2}=y,
\]
\[
  E_{k}=\frac{1}{2}mv_{0}''^{2}+mgy,
\]
联立解得
\[
  x^{2}+4y^{2}-8\mu ly=0\qquad(0<x\le2\mu l).
\]

(3) 由动能定理得
\[
  mgh-\mu mgl=\frac{1}{2}mv_{1}^{2}-0.
\]
碰撞时（$A$、$B$ 弹性碰撞，$A$ 碰后反弹）有
\[
  mv_{1}=-mv_{2}+\lambda mv_{B},
\]
\[
  \frac{1}{2}mv_{1}^{2}=\frac{1}{2}mv_{2}^{2}+\frac{1}{2}\lambda mv_{B}^{2}.
\]
设 $A$ 碰撞后又回到 $O$ 点的速度为 $v_{3}$，有
\[
  -\mu mg\cdot 2l=\frac{1}{2}mv_{3}^{2}-\frac{1}{2}mv_{2}^{2}.
\]
$A$ 落到弧形轨道上，故
\[
  v_{3}\le v_{0}',
\]
$A$ 落到 $B$ 落点右侧，故
\[
  v_{3}>v_{B}.
\]
联立解得
\[
  \frac{3\lambda-1}{\lambda-3}\,\mu l<h\le\frac{4(\lambda^{2}+\lambda+1)}{(\lambda-1)^{2}}\,\mu l.
\]
\end{answer}

\begin{answer}
\gksolhead{15}{选考题（选修 3-3）}
\gkanswer{答案 (1) BDE；(2)(i) $T_{2}=\SI{297}{\kelvin}$；(2)(ii) $T_{\max}=\SI{309}{\kelvin}$}
(1) 活塞 $S_{2}$ 没有移动，所以外力 $F$ 做功为零，选项 A 错误。由于环境温度不变，汽缸导热，汽缸内理想气体温度不变，分子平均动能不变，所以理想气体内能不变，选项 B 正确、C 错误。整个过程中，因为 $m$ 是不断增加的，压力做功小于 $p_{0}S_{1}h+mgh$，根据热力学第一定律可知，理想气体向外界释放的热量小于 $p_{0}S_{1}h+mgh$，选项 D 正确。选右侧活塞为研究对象，根据平衡条件有
\[
  F+p_{0}S_{2}=pS_{2};
\]
选左侧活塞为研究对象，根据平衡条件有
\[
  mg+p_{0}S_{1}=pS_{1},
\]
联立解得
\[
  F=\frac{mgS_{2}}{S_{1}},
\]
选项 E 正确。

(2)(i) 对铁块受力分析，根据平衡条件有
\[
  F_{T}=m_{2}g-m_{0}g,
\]
其中 $m_{0}$ 为天平示数。左侧理想气体初状态：
\[
  p_{1}=p_{0}+\frac{m_{1}g-F_{T1}}{S},\qquad T_{1}=\SI{300}{\kelvin};
\]
末状态：
\[
  p_{2}=p_{0}+\frac{m_{1}g-F_{T2}}{S},\qquad T_{2}\ \text{待求}.
\]
根据等容变化，得
\[
  \frac{p_{1}}{T_{1}}=\frac{p_{2}}{T_{2}},
\]
联立解得 $T_{2}=\SI{297}{\kelvin}$。

(2)(ii) 当环境温度为 $T_{\max}$ 时，天平读数为 \SI{1200}{\gram}，此时 $F_{T}=0$，
\[
  p_{3}=p_{0}+\frac{m_{1}g}{S},
\]
根据等容变化，得
\[
  \frac{p_{1}}{T_{1}}=\frac{p_{3}}{T_{\max}},
\]
联立解得 $T_{\max}=\SI{309}{\kelvin}$。
\end{answer}
